\input{preambulo}

\begin{document}

\maketitle

\begin{abstract} % 150-250 words; problem, solution, objectives, expected outcomes.
    \textbf{COMPLETAMENTE TEMPORÁRIO}:
    % Problem
    Muitos produtos de e-commerce possuem múltiplos atributos (ex.: tamanho, cor, material) que são essenciais para que os clientes tomem decisões de compra informadas. Entretanto, extrair e estruturar essas informações a partir de descrições de produtos não estruturadas é uma tarefa desafiadora.
    % Solution
    Nós propomos o JSONLLM, uma nova abordagem que aproveita os grandes modelos de linguagem (LLMs) para extrair e estruturar atributos de produtos em formato JSON. Ao utilizar as avançadas capacidades de compreensão de linguagem natural dos Grandes Modelos de Linguagem, nosso método visa identificar e categorizar com precisão os atributos dos produtos a partir de descrições diversas e complexas.
    % Objectives
    Os objetivos principais deste projeto incluem: (1) desenvolver um pipeline que integra LLMs para a extração de atributos de produtos; (2) criar um esquema padronizado para representar os atributos extraídos; (3) avaliar a precisão e eficiência do modelo proposto em comparação com métodos tradicionais de extração de informações.
    % Expected outcomes
    Esperamos que este sistema melhore significativamente a precisão da extração de atributos de produtos, facilitando a criação de catálogos de produtos mais ricos e detalhados. Além disso, a estruturação em JSON permitirá uma integração mais fácil com sistemas de e-commerce existentes, melhorando a experiência do usuário final.
\end{abstract}

\section{Introdução e Motivação} % Introduction and Motivation: Background, problem statement, significance, motivation.

\section{Objetivos} % Project Objectives: SMART objectives (clear, specific, measurable, achievable, relevant, and time-bound).

\section{Metodologia} % Proposed Methodology: Data collection/preparation, model/architecture selection, implementation details, experimental design, evaluation plan.

\section{Resultados Esperados e Contribuições} % Expected Outcomes and Contributions: Anticipated results, contributions to LLM field, potential applications.

\section{Cronograma} % Timeline and Milestones: 1. Literature Review; 2. Data Preparation; 3. Model Implementation; 4. Experimentation; 5. Report Draft; 6. Final Submission.

\section{Referências} % References: Cite relevant literature with consistent formatting.

\section{Desafios e Estratégias de Mitigação} % Potential Challenges and Mitigation Strategies: Potential challenges and mitigation strategies.

\end{document}
